\section{第一阶段深讲：设备、总线、DMA 与中断的内核契约}

设备是操作系统最容易讲浅的部分。把驱动说成“读写寄存器”，就会漏掉总线枚举、资源分配、DMA 地址、队列所有权、中断亲和性、电源状态、错误恢复和安全隔离。本节把设备看作一个和 CPU 并行运行的状态机：内核通过寄存器和共享内存给它任务，它通过中断和 completion 把结果交回来。

\subsection{设备不是函数调用对象}

函数调用有明确进入和返回。设备请求不是这样。内核提交请求后，设备可能立即完成，也可能稍后完成，也可能超时，也可能在复位后返回错误。请求对象必须在这段时间内保持身份稳定。

\begin{lstlisting}
device request lifecycle:
  allocated
  initialized
  mapped for DMA
  visible in submission queue
  owned by device
  completed or timed out
  unmapped
  reported to upper layer
  released
\end{lstlisting}

这个生命周期解释了为什么驱动代码到处有引用计数、状态位、超时定时器和完成回调。没有它们，内核无法判断一个 buffer 现在归 CPU 还是设备，无法处理迟到 completion，也无法在进程退出时安全取消请求。

\subsection{总线枚举决定设备怎样被发现}

不同总线的发现方式差别很大。PCIe 设备有配置空间，可以枚举 vendor、device、class 和 BAR；USB 设备有描述符和接口；I2C/SPI 这类板载设备常不能完整自描述，在本书 x86-64 主线里通常需要 ACPI 提供地址、IRQ 和电源资源。OS 的 driver model 就是为这些差异建立统一框架。

\begin{longtable}{p{0.18\linewidth}p{0.34\linewidth}p{0.34\linewidth}}
\toprule
总线 & 发现方式 & 驱动关心的资源 \\
\midrule
PCIe & 配置空间扫描 & BAR、MSI-X、DMA mask、AER \\
USB & 设备和接口描述符 & endpoint、URB、电源、热插拔 \\
I2C & 固件描述地址和类型 & 设备地址、中断 GPIO、时钟 \\
SPI & 固件描述片选和模式 & 片选、频率、mode、DMA 能力 \\
platform bus & ACPI 描述 & MMIO、IRQ、clock、reset、regulator \\
\bottomrule
\end{longtable}

一个驱动的 probe 函数不应假设资源永远存在。它要检查 MMIO 映射是否成功，中断是否申请成功，DMA mask 是否能覆盖需要地址，设备固件是否加载成功，电源域是否已经打开。probe 失败要逆序释放已经申请的资源。

\subsection{MMIO 寄存器是设备协议，不是普通变量}

MMIO 读写会产生副作用。读状态寄存器可能清 pending 位；写控制寄存器可能启动设备；写 doorbell 可能让设备读取共享队列。编译器和 CPU 对普通内存的优化不能随便套到 MMIO 上，所以内核使用专门的读写接口。

\begin{lstlisting}
submit queue entry:
  sq[tail].addr = dma_addr
  sq[tail].len = len
  sq[tail].cmd = READ
  dma_wmb()
  writel(new_tail, device_doorbell)
\end{lstlisting}

这个顺序有协议含义。设备看到 doorbell 时，必须已经能读到完整描述符。若把 \code{writel} 当普通 store，或者缺少 DMA 写屏障，设备可能读到旧内容。驱动 bug 的随机性，很多来自这种硬件可见顺序。

\subsection{DMA 地址不是用户虚拟地址}

设备不能拿用户虚拟地址直接 DMA。用户虚拟地址只对某个进程和 CPU MMU 有意义。设备看到的是总线地址或 IOVA。内核必须把参与 DMA 的页 pin 住或使用内核 buffer，再建立设备可用地址映射。

\begin{lstlisting}
prepare DMA:
  choose buffer pages
  ensure pages remain resident
  map pages for device direction
  get dma_addr visible to device
  program descriptor with dma_addr
\end{lstlisting}

DMA 方向也很重要。ToDevice 表示 CPU 写好数据，设备读；FromDevice 表示设备写，CPU 之后读；Bidirectional 两者都有。方向决定 cache clean/invalidate 和 IOMMU 权限。方向标错可能表现为旧数据、数据损坏或 IOMMU fault。

\subsection{IOMMU 给设备建立地址空间}

没有 IOMMU，设备拿到总线地址后可能访问很大物理范围。IOMMU 让内核为设备建立受限映射：某个 IOVA 只对应被授权的物理页，方向和权限受控。这样设备或驱动出错时，不至于随意覆盖内核内存。

\begin{longtable}{p{0.2\linewidth}p{0.34\linewidth}p{0.34\linewidth}}
\toprule
机制 & 作用 & 代价 \\
\midrule
IOVA & 给设备看的地址 & 需要映射管理和失效 \\
IOTLB & 缓存设备侧翻译 & unmap 后要处理失效 \\
domain & 一组设备的隔离空间 & 设备分组和直通复杂 \\
DMA remapping & 限制设备可访问页 & 增加映射开销 \\
\bottomrule
\end{longtable}

虚拟化和设备直通尤其依赖 IOMMU。把设备交给虚拟机时，宿主必须确保设备只能访问该虚拟机的内存，而不能 DMA 到宿主或其他虚拟机。

\subsection{中断处理要分快慢两层}

硬件中断入口运行在受限上下文里，不能做太多慢操作。典型设计是 top half 尽快确认设备、取走必要状态、屏蔽或重新使能中断，然后把耗时工作推到 softirq、tasklet、workqueue、threaded IRQ 或 NAPI poll 中。

\begin{lstlisting}
interrupt top half:
  read device interrupt status
  acknowledge or mask interrupt
  record that work is pending
  schedule bottom half
  return from interrupt

bottom half:
  process completion queue
  unmap DMA buffers
  wake waiters
  submit follow-up work if needed
\end{lstlisting}

分层的原因是延迟和可重入。中断关闭太久会影响定时器和其他设备；在硬中断里睡眠会破坏调度语义；做大量包处理会造成中断风暴。网络 NAPI 就是把高频中断转换为预算受控的轮询，减少每包中断成本。

\subsection{MSI-X 和中断亲和性连接到多核性能}

现代高速设备常有多个队列，每个队列配一个或多个 MSI-X vector。OS 可以把不同队列的中断分配到不同 CPU，让网络包或存储完成更接近处理它们的线程和内存。

\begin{lstlisting}
multi-queue device:
  queue 0 interrupt -> CPU 0
  queue 1 interrupt -> CPU 1
  queue 2 interrupt -> CPU 2
  queue 3 interrupt -> CPU 3
\end{lstlisting}

若所有中断都落到 CPU0，CPU0 会成为瓶颈，其他 CPU 等待数据。若中断在 CPU 和内存 node 之间乱跑，会增加 cache 和 NUMA 成本。设备队列、IRQ affinity、RSS/RPS、块层 tag 和用户线程绑核，是同一个性能问题的不同层次。

\subsection{设备队列把硬件并行度暴露给内核}

NVMe、网卡、GPU 等设备都通过队列表达并行能力。队列深度越大，设备能同时处理的请求越多，但排队过深也会增加尾延迟。内核要在吞吐和延迟之间控制 in-flight 请求数量。

\begin{longtable}{p{0.2\linewidth}p{0.34\linewidth}p{0.34\linewidth}}
\toprule
队列状态 & 含义 & OS 策略 \\
\midrule
submission queue & CPU 写入设备要做的命令 & 合并、排序、限流、doorbell \\
completion queue & 设备写入已完成命令 & phase tag、轮询、中断 \\
queue depth & 同时在飞请求数量 & 吞吐和尾延迟折中 \\
tag/request id & 请求身份 & 乱序完成、取消、超时处理 \\
\bottomrule
\end{longtable}

设备完成可以乱序返回，所以 request id 必须稳定。若复用 tag 过早，旧 completion 可能错误匹配新请求。NVMe completion 的 phase tag 正是为了解决环形队列槽位复用问题。

\subsection{电源管理是设备状态机的一部分}

设备空闲时关闭时钟、电源域或进入低功耗状态可以省电，但 resume 后寄存器和队列状态可能丢失。驱动必须知道设备从 suspended 到 active 需要哪些步骤。

\begin{lstlisting}
runtime suspend:
  stop new requests
  drain or cancel in-flight requests
  save registers if required
  disable interrupts
  gate clocks and power

runtime resume:
  enable power and clocks
  restore registers and queues
  reenable interrupts
  accept requests again
\end{lstlisting}

若 suspend 时仍有 DMA，设备可能写入已经被释放的 buffer。若 resume 后忘记恢复 doorbell 或队列地址，提交请求会超时。电源管理 bug 常表现为“第一次正常，睡眠唤醒后失败”。

\subsection{错误恢复比正常路径更能检验驱动质量}

设备可能超时、返回错误 completion、发生 PCIe AER、被用户拔出、固件崩溃或电源丢失。驱动必须把错误传播到上层，而不是让任务永久睡眠。

\begin{lstlisting}
timeout handler:
  mark request timed out
  stop affected queue
  attempt device reset
  complete request with error or retry
  wake waiters
  unfreeze queue if recovery succeeded
\end{lstlisting}

错误恢复要特别防止双完成。一个请求超时后，驱动可能已经向上层返回错误；如果设备随后给出迟到 completion，驱动必须识别它属于旧请求，不能再次完成或释放同一个对象。引用计数和 generation id 常用于解决这类问题。

\subsection{设备安全边界}

设备驱动运行在内核里，攻击面很大。用户态通过 \code{ioctl}、mmap、网络包、USB 描述符或 GPU 命令间接影响驱动。所有来自设备和用户的数据都必须验证。IOMMU 限制 DMA，权限系统限制谁能打开设备文件，seccomp/LSM 可以继续收紧接口。

\begin{longtable}{p{0.2\linewidth}p{0.34\linewidth}p{0.34\linewidth}}
\toprule
攻击面 & 风险 & 防线 \\
\midrule
用户 ioctl 参数 & 越界、类型混淆、UAF & 参数复制、长度检查、对象引用 \\
设备 DMA & 覆盖任意内存 & IOMMU、DMA API、方向权限 \\
恶意 USB 设备 & 描述符畸形、热插拔竞态 & 描述符验证、授权策略 \\
网络包 & 协议解析漏洞 & 长度检查、分层解析、fuzz \\
GPU 命令 & 访问其他进程显存 & IOMMU、命令验证、上下文隔离 \\
\bottomrule
\end{longtable}

这就是为什么“驱动只是硬件适配层”是不准确的。驱动同时是性能路径、权限边界、错误恢复层和硬件状态机管理者。

\subsection{设备契约的最低验收}

第一阶段应能解释：为什么 MMIO 访问需要专门接口；为什么 DMA buffer 在 completion 前不能释放；为什么 IOMMU 不是可有可无的性能组件；为什么中断要分 top half 和 bottom half；为什么多队列设备要考虑中断亲和性；为什么驱动 reset 必须处理所有 in-flight 请求。

能回答这些问题，后面学习块层、网络栈、文件系统和 io\_uring 才不会停留在 API 表面。
